\documentclass[11pt,a4paper]{article}

\usepackage{fontspec}
\usepackage[croatian]{babel}
\usepackage[a4paper,margin=2.2cm,top=2.4cm,bottom=2.4cm]{geometry}
\usepackage{xcolor}
\usepackage{graphicx}
\usepackage{array,tabularx,booktabs,colortbl}
\usepackage{enumitem}
\usepackage{fancyvrb}
\usepackage[most]{tcolorbox}
\usepackage{titlesec}
\usepackage{fancyhdr}
\usepackage{lastpage}
\usepackage{needspace}
\usepackage[hidelinks]{hyperref}

\setmainfont{IBMPlexSans}[
  Extension=.otf, UprightFont=*-Regular, ItalicFont=*-Italic,
  BoldFont=*-SemiBold, BoldItalicFont=*-SemiBoldItalic]
\setsansfont{IBMPlexSans}[
  Extension=.otf, UprightFont=*-Regular, ItalicFont=*-Italic,
  BoldFont=*-SemiBold, BoldItalicFont=*-SemiBoldItalic]
\setmonofont{IBMPlexMono}[
  Extension=.otf, UprightFont=*-Regular, ItalicFont=*-Italic,
  BoldFont=*-SemiBold, BoldItalicFont=*-SemiBoldItalic, Scale=0.9]

\definecolor{ink}{HTML}{0C4685}
\definecolor{ink2}{HTML}{6E7BAA}
\definecolor{tint}{HTML}{EAEEF6}
\definecolor{tekst}{HTML}{1A222E}
\definecolor{muted}{HTML}{6B7484}
\definecolor{pozadina}{HTML}{F6F7F9}
\definecolor{linije}{HTML}{D9DEE7}

\color{tekst}
\setlength{\parindent}{0pt}
\setlength{\parskip}{0.55em}
\linespread{1.12}
\hypersetup{colorlinks=true,urlcolor=ink,linkcolor=ink}

\titleformat{\section}{\Large\bfseries\color{ink}}{}{0pt}{}
\titleformat{\subsection}{\large\bfseries\color{ink}}{}{0pt}{}
\titlespacing*{\section}{0pt}{1.4em}{0.4em}
\titlespacing*{\subsection}{0pt}{1.1em}{0.3em}

\pagestyle{fancy}
\fancyhf{}
\renewcommand{\headrulewidth}{0pt}
\fancyhead[L]{\footnotesize\color{muted}Programiranje za web - prelazak na novo predavanje}
\fancyfoot[R]{\footnotesize\color{muted}Stranica \thepage\ od \pageref*{LastPage}}

\newcommand{\code}[1]{{\ttfamily\small #1}}
\newcommand{\provjera}{\textbf{\color{ink}Provjera.}~}
\newcolumntype{L}{>{\raggedright\arraybackslash}X}
\newcolumntype{P}[1]{>{\raggedright\arraybackslash}p{#1}}
\renewcommand{\arraystretch}{1.35}

\newenvironment{kod}
  {\VerbatimEnvironment
   \begin{tcolorbox}[colback=pozadina,colframe=linije,boxrule=0.4pt,arc=2pt,
     left=6pt,right=6pt,top=3pt,bottom=3pt]
   \begin{Verbatim}[fontsize=\small]}
  {\end{Verbatim}\end{tcolorbox}}

\newenvironment{ispis}
  {\VerbatimEnvironment
   \begin{tcolorbox}[colback=pozadina,colframe=linije,boxrule=0.4pt,arc=2pt,
     left=6pt,right=6pt,top=3pt,bottom=3pt]
   \begin{Verbatim}[fontsize=\footnotesize]}
  {\end{Verbatim}\end{tcolorbox}}

% Ispis širi od stranice: spremi s SaveVerbatim, pa \siroko{naziv}
\newcommand{\siroko}[1]{%
  \begin{tcolorbox}[colback=pozadina,colframe=linije,boxrule=0.4pt,arc=2pt,
    left=6pt,right=6pt,top=3pt,bottom=3pt]
  \resizebox{\linewidth}{!}{\BUseVerbatim{#1}}
  \end{tcolorbox}}

% Snimka zaslona iz mape slike/; dok je nema, okvir s nazivom datoteke
\newcommand{\slika}[3][0.85]{%
  \begin{center}\begin{minipage}{\linewidth}\centering
  \IfFileExists{slike/#2}%
    {\fcolorbox{linije}{white}{\includegraphics[width=#1\linewidth]{slike/#2}}}%
    {\fcolorbox{linije}{pozadina}{\parbox[c][3.2cm][c]{0.8\linewidth}{\centering
      \color{muted}\small snimka zaslona: \code{slike/#2}}}}%
  \par\vspace{0.2em}{\small\color{muted}#3}
  \end{minipage}\end{center}}

\newcommand{\zaglavlje}{\rowcolor{tint}}

\begin{document}

{\raggedright\hyphenpenalty=10000\color{ink}\fontsize{22}{26}\selectfont\bfseries Programiranje za web - prelazak na novo predavanje\par}
\vspace{0.3em}
{\color{muted}18. listopada 2026. \textperiodcentered{} doc. dr. sc. Tomislav Petković\par}
\vspace{0.4em}

\section{Uvod}

Svako predavanje ima svoj tag u repozitoriju projekta: \code{predavanje-01},
\code{predavanje-02} i tako dalje. Kod se između dva taga mijenja, a s njim i shema baze,
popis paketa i konfiguracija. Zbog toga okruženje iz prethodnog predavanja može biti
nedostatno za rad na trenutnom predavanju.

Ova vježba prolazi cijeli prelazak s jednog predavanja na drugo, korak po korak. Iza
svakog koraka stoji \textbf{provjera}: naredba i ispis koji morate vidjeti prije nego
krenete dalje. Ako ispis nije isti, ne nastavljajte; pogledajte odjeljak
\hyperref[sec:problemi]{\emph{Česti problemi}} na kraju.

Primjeri ispisa u nastavku su za prelazak na \code{predavanje-03}. Za svako sljedeće
predavanje postupak je isti, a razlikuju se broj migracija, sadržaj baze i rute koje je
potrebno provjeriti.

\subsection{Dva terminala}

U VS Code-u otvorite mapu repozitorija (\emph{File → Open Folder}) i u njoj dva terminala
(gumb \textbf{+} u panelu terminala):

\begin{center}\small
\begin{tabularx}{\linewidth}{P{2.2cm}P{3.6cm}L}
\arrayrulecolor{linije}\toprule
\zaglavlje \textbf{Terminal} & \textbf{Mapa} & \textbf{Za što}\\
\midrule
prvi & korijen repozitorija & \code{git}, \code{docker compose}\\
drugi & \code{api} & virtualno okruženje, \code{alembic}, \code{python -m app.seed}, \code{uvicorn}\\
\bottomrule
\end{tabularx}
\end{center}

Naredbe \code{docker compose} traže datoteku \code{docker-compose.yml}, koja stoji
u korijenu. Pokrenute iz mape \code{api} javljaju
\code{no configuration file provided: not found}.

\section{0. Kraj prethodnog predavanja}

Ako ste prošli put završili prema koraku 3 ove vježbe, baza je već obrisana i ovaj korak
preskočite. Inače, u prvom terminalu:

\begin{kod}
docker compose down -v
\end{kod}

Zašto se baza briše opisuje korak 2b.

\section{1. Kod novog predavanja}

\subsection{1a. Čisto radno stablo}

\begin{kod}
git status
\end{kod}

\provjera Ispis završava retkom:

\begin{ispis}
nothing to commit, working tree clean
\end{ispis}

Ako \code{git status} navodi izmijenjene datoteke (\code{modified:}) ili nove datoteke
(\code{Untracked files:}), radno stablo nije čisto. Izmjene obrišite:

\begin{kod}
git restore .       # izmijenjene datoteke
git clean -fd       # nove datoteke
\end{kod}

Ponovite \code{git status} dok ne dobijete \code{working tree clean}. Tko izmjene želi
sačuvati, može ih prije brisanja spremiti na svoju granu.

\subsection{1b. Prelazak na tag}

\begin{kod}
git checkout predavanje-03
git status
\end{kod}

\provjera Ispis zadnje naredbe:

\begin{ispis}
HEAD detached at predavanje-03
nothing to commit, working tree clean
\end{ispis}

\emph{Detached HEAD} znači da niste ni na jednoj grani, nego na točno određenom
\emph{commitu}. Za praćenje predavanja to je ispravno stanje.

\subsection{1c. Što checkout ne mijenja}

Neke datoteke git namjerno ne prati. One ostaju na disku kakve jesu, bez obzira na tag:

\begin{kod}
git status --ignored --short
\end{kod}

\begin{ispis}
!! .env
!! api/.venv/
!! api/app/__pycache__/
...
\end{ispis}

To su upravo \code{.env} s konfiguracijom i \code{api/.venv} s instaliranim paketima. Kod
je novi, a one su iz prošlog tjedna. Zbog toga je potrebno provjeriti korake 2a i 2c.

\section{2. Pokretanje prema README-u}

Otvorite \code{README.md} novog predavanja i pronađite odjeljak \emph{Pokretanje}. Koraci
u nastavku slijede ga redom, uz provjeru iza svakoga.

\subsection{2a. Konfiguracija}

Predložak \code{.env.example} mijenja se između predavanja, pa se \code{.env} pri svakom
prelasku kopira iznova. U prvom terminalu:

\begin{kod}
cp .env.example .env        # PowerShell: Copy-Item .env.example .env
\end{kod}

Naredba prepisuje staru \code{.env}. Ako kopiranje javi pogrešku, obrišite staru
\code{.env} i ponovite naredbu.

\subsection{2b. Baza}

\textbf{Prvo provjerite da nijedna baza ne postoji.} Pokrenite Docker Desktop i pričekajte
\emph{Engine running}. U prvom terminalu:

\begin{kod}
docker ps -a
docker volume ls
\end{kod}

\code{docker ps -a} prikazuje sve kontejnere na računalu: one koji rade, one koji su
ugašeni i one iz drugih projekata.

\provjera Obje naredbe smiju ispisati samo zaglavlje:

\begin{ispis}
CONTAINER ID   IMAGE     COMMAND   CREATED   STATUS    PORTS     NAMES
DRIVER    VOLUME NAME
\end{ispis}

Primjer ispisa kad još radi baza iz uputa za pripremu okruženja (\code{pzw-provjera}),
a postoji i baza ovog projekta:

\begin{SaveVerbatim}{dockerps}
CONTAINER ID   IMAGE         COMMAND                  CREATED         STATUS        PORTS                                         NAMES
bb50f3845718   postgres:16   "docker-entrypoint.s…"   1 second ago    Created                                                     situnizd-pzw-projekt-db-1
4df8d4a6dce7   postgres:17   "docker-entrypoint.s…"   2 seconds ago   Up 1 second   0.0.0.0:5432->5432/tcp, [::]:5432->5432/tcp   pzw-provjera-db-1

DRIVER    VOLUME NAME
local     613f4956b73f98316fc6b44486aef6a4d393870a8c030b16a6087fc59953a570
local     situnizd-pzw-projekt_pgdata
\end{SaveVerbatim}
\siroko{dockerps}

Oba kontejnera traže isti port 5432. Dok \code{pzw-provjera-db-1} radi, baza ovog projekta
ne može se pokrenuti i ostaje u stanju \code{Created}, a \code{docker compose up} javlja
\code{Bind for 0.0.0.0:5432 failed: port is already allocated}.

\textbf{Obrišite svaki kontejner} po imenu iz stupca \code{NAMES}:

\begin{kod}
docker rm -f -v pzw-provjera-db-1
docker rm -f -v situnizd-pzw-projekt-db-1
\end{kod}

Zastavica \code{-f} kontejner koji radi prvo zaustavi, pa ga obriše. Zastavica \code{-v} uz
kontejner obriše i njegov volumen bez imena; takav je volumen s nazivom od dugog niza
znakova u gornjem ispisu, koji je stvorio \code{pzw-provjera}. Volumen s imenom ostaje, pa
ga \textbf{obrišite} po imenu iz stupca \code{VOLUME NAME}:

\begin{kod}
docker volume rm situnizd-pzw-projekt_pgdata
\end{kod}

Volumen se ne može obrisati dok postoji kontejner koji ga koristi, pa se kontejneri brišu
prvi. Ako Docker koristite i za druge kolegije, brišite samo kontejnere i volumene ovog
kolegija.

Ponovite \code{docker ps -a} i \code{docker volume ls} dok obje ne ispišu samo zaglavlje.

Baza se briše zato što pamti stanje starog koda: tablice, stupce i broj zadnje migracije.
Novo predavanje može tablicu preimenovati ili joj dodati stupac, a stari podatci tada ne
odgovaraju novom kodu. Najmanje problema stvara prazna baza koju gradi kod novog predavanja.

\textbf{Pokretanje baze:}

\begin{kod}
docker compose up -d db
docker compose ps
\end{kod}

\provjera U ispisu \code{up} mora pisati da je volumen stvoren (\code{Created}), a ne
zatečen:

\begin{ispis}
 Volume situnizd-pzw-projekt_pgdata Created
 Network situnizd-pzw-projekt_default Created
 Container situnizd-pzw-projekt-db-1 Started
\end{ispis}

Druga naredba prikazuje red s \code{Up} i \code{0.0.0.0:5432->5432/tcp}:

\begin{SaveVerbatim}{psispis}
NAME                        IMAGE         COMMAND                  SERVICE   CREATED         STATUS         PORTS
situnizd-pzw-projekt-db-1   postgres:16   "docker-entrypoint.s…"   db        5 seconds ago   Up 4 seconds   0.0.0.0:5432->5432/tcp, [::]:5432->5432/tcp
\end{SaveVerbatim}
\siroko{psispis}

\textbf{Provjera da je baza prazna, naredbom.} Klijent \code{psql} već postoji
u kontejneru s bazom:

\begin{kod}
docker compose exec db psql -U pl_user -d pl_reg -c "\dt"
\end{kod}

\begin{ispis}
Did not find any relations.
\end{ispis}

Baza \code{pl\_reg} postoji, ali u njoj nema nijedne tablice.

\textbf{Provjera da je baza prazna, ekstenzijom Database.} Na lijevoj traci VS Code-a
otvorite \emph{Database} i dodajte novu vezu tipa PostgreSQL:

\begin{center}\small
\begin{tabularx}{\linewidth}{P{4cm}L}
\arrayrulecolor{linije}\toprule
\zaglavlje \textbf{Podatak} & \textbf{Vrijednost}\\
\midrule
Host & \code{127.0.0.1}\\
Port & \code{5432}\\
Username & \code{pl\_user}\\
Password & \code{pl\_pass}\\
Database & \code{pl\_reg}\\
\bottomrule
\end{tabularx}
\end{center}

Vrijednosti su iste kao u datoteci \code{docker-compose.yml} i u varijabli
\code{DATABASE\_URL} u \code{.env}. Veza ostaje zapamćena, pa je sljedeći tjedan ne treba
ponovno upisivati.

Nakon spajanja raširite \code{127.0.0.1@5432 → pl\_reg → public → Tables}. Ondje piše
\emph{No tables found.}

\slika[0.42]{baza-prazna.png}{Database: prazna baza}

\subsection{2c. Virtualno okruženje i paketi}

U drugom terminalu, u mapi \code{api}:

\begin{kod}
cd api
\end{kod}

Ako mapa \code{.venv} ne postoji, napravite je s \code{uv venv}. Ako postoji, ne pokrećite
\code{uv venv} ponovno, jer javlja \code{A virtual environment already exists at: .venv}.

Aktivacija:

\begin{kod}
.venv\Scripts\Activate.ps1      # Windows (PowerShell)
source .venv/bin/activate       # macOS / Linux
\end{kod}

\provjera Ispred prompta mora pisati \code{(api)} ili \code{(.venv)}. Sigurnija provjera
je pitati Python gdje je:

\begin{kod}
python -c "import sys; print(sys.prefix)"
\end{kod}

\begin{ispis}
...\situnizd-pzw-projekt\api\.venv
\end{ispis}

Ako ispis ne završava s \code{api\textbackslash.venv} (na macOS-u i Linuxu
\code{api/.venv}), paketi bi se instalirali u sistemski Python. Aktivirajte okruženje
i ponovite provjeru.

\textbf{Instalacija paketa.} Popis u \code{requirements.txt} mijenja se između predavanja,
pa se ova naredba pokreće pri \textbf{svakom} prelasku, i kad \code{.venv} već postoji:

\begin{kod}
uv pip install -r requirements.txt
\end{kod}

\provjera Paketi koji su već instalirani se preskaču, a novi se dodaju. Primjer pri
prelasku s predavanja 2 na predavanje 3:

\begin{ispis}
Resolved 39 packages in 1.49s
Installed 11 packages in 466ms
 + httpx==0.28.1
 + pyjwt==2.15.1
 + pytest==9.1.1
 ...
\end{ispis}

Bez ovoga koraka aplikacija pada pri pokretanju s
\code{ModuleNotFoundError: No module named 'jwt'}.

\subsection{2d. Migracije i početni podatci}

U drugom terminalu, s aktivnim virtualnim okruženjem:

\begin{kod}
alembic upgrade head
\end{kod}

\begin{SaveVerbatim}{alembicispis}
INFO  [alembic.runtime.migration] Context impl PostgresqlImpl.
INFO  [alembic.runtime.migration] Will assume transactional DDL.
INFO  [alembic.runtime.migration] Running upgrade  -> 001, init clubs and users
INFO  [alembic.runtime.migration] Running upgrade 001 -> 002, rename user email to username, add is_active
\end{SaveVerbatim}
\siroko{alembicispis}

Svaki redak \code{Running upgrade} je jedna migracija iz mape \code{api/alembic/versions}.
U praznoj bazi primjenjuju se sve, od prve. Na kojoj je migraciji baza sada:

\begin{kod}
alembic current
\end{kod}

\begin{ispis}
002 (head)
\end{ispis}

\code{(head)} znači da je baza na zadnjoj migraciji koju kod poznaje.

Početni podatci:

\begin{kod}
python -m app.seed
\end{kod}

\begin{ispis}
Seed gotov.
  organizator: admin@plreg.hr / admin123
  klub:        behemot@plreg.hr / behemot123
\end{ispis}

\textbf{Provjera sadržaja baze, naredbom.} U prvom terminalu:

\begin{kod}
docker compose exec db psql -U pl_user -d pl_reg -c "\dt"
\end{kod}

\begin{ispis}
             List of relations
 Schema |      Name       | Type  |  Owner
--------+-----------------+-------+---------
 public | alembic_version | table | pl_user
 public | clubs           | table | pl_user
 public | users           | table | pl_user
(3 rows)
\end{ispis}

Tablicu \code{alembic\_version} stvara Alembic i u njoj pamti broj zadnje primijenjene
migracije. Ostale tablice odgovaraju modelima u \code{api/app/models}.

\begin{SaveVerbatim}{selectkod}
docker compose exec db psql -U pl_user -d pl_reg -c "SELECT id, username, role, club_id, is_active FROM users;"
\end{SaveVerbatim}
\siroko{selectkod}

\begin{ispis}
 id |     username     | role  | club_id | is_active
----+------------------+-------+---------+-----------
  1 | admin@plreg.hr   | admin |         | t
  2 | behemot@plreg.hr | club  |       1 | t
(2 rows)
\end{ispis}

Stupac \code{password\_hash} je namjerno izostavljen iz upita; lozinke su u bazi spremljene
kao \emph{hash}, ne kao tekst.

\textbf{Provjera sadržaja baze, ekstenzijom Database.} Osvježite popis tablica (ikona
osvježavanja pokraj \emph{Tables}). Pod \emph{Tables} su sada \code{alembic\_version},
\code{clubs} i \code{users}. Klik na tablicu otvara njezine retke; na slici je tablica
\code{clubs} s jednim klubom.

\slika[0.95]{baza-seed.png}{Database: tablice nakon migracija i retci tablice \code{clubs}}

Za predavanje 3 sadržaj baze je:

\begin{center}\small
\begin{tabularx}{\linewidth}{P{3.6cm}P{1.8cm}L}
\arrayrulecolor{linije}\toprule
\zaglavlje \textbf{Tablica} & \textbf{Redaka} & \textbf{Sadržaj}\\
\midrule
\code{alembic\_version} & 1 & \code{002}\\
\code{clubs} & 1 & PLK Behemot, Zadar\\
\code{users} & 2 & organizator \code{admin@plreg.hr}, račun kluba \code{behemot@plreg.hr}\\
\bottomrule
\end{tabularx}
\end{center}

\subsection{2e. Pokretanje aplikacije}

U drugom terminalu:

\begin{kod}
uvicorn app.main:app --reload
\end{kod}

\provjera Ispis završava retkom:

\begin{ispis}
INFO:     Application startup complete.
\end{ispis}

Terminal ostavite otvoren dok aplikacija radi. U pregledniku otvorite:

\begin{center}\small
\begin{tabularx}{\linewidth}{P{6cm}L}
\arrayrulecolor{linije}\toprule
\zaglavlje \textbf{Adresa} & \textbf{Očekivano}\\
\midrule
\url{http://127.0.0.1:8000/health} & \code{\{"status":"ok"\}}\\
\url{http://127.0.0.1:8000/docs} & Swagger UI s popisom ruta\\
\bottomrule
\end{tabularx}
\end{center}

U Swaggeru otvorite rutu \code{GET /health/}, pritisnite \emph{Try it out}, zatim
\emph{Execute}. Ispod se prikazuju statusni kod \code{200} i tijelo odgovora
\code{\{"status": "ok"\}}.

\slika[0.8]{swagger-health.png}{Swagger: odgovor na \code{GET /health}}

U terminalu s aplikacijom svaki zahtjev ostavlja redak sa statusnim kodom:

\begin{ispis}
INFO:     127.0.0.1:50852 - "GET /health/ HTTP/1.1" 200 OK
\end{ispis}

\section{3. Gašenje i provjera da je sve prazno}

Zaustavite aplikaciju s \textbf{Ctrl+C} u drugom terminalu i izađite iz virtualnog
okruženja:

\begin{kod}
deactivate
\end{kod}

U prvom terminalu ugasite bazu i obrišite njezin volumen:

\begin{kod}
docker compose down -v
\end{kod}

\begin{ispis}
 Container situnizd-pzw-projekt-db-1 Removed
 Volume situnizd-pzw-projekt_pgdata Removed
 Network situnizd-pzw-projekt_default Removed
\end{ispis}

Zastavica \code{-v} briše i volumen. Bez nje bi \code{down} uklonio kontejner, a podatci bi
ostali i sljedeći tjedan dočekali kod novog predavanja. Sve što ste unijeli kroz Swagger
time nestaje; početne podatke vraća \code{python -m app.seed}.

\textbf{Provjera naredbama:}

\begin{kod}
docker ps -a
docker volume ls
git status
\end{kod}

Prve dvije ispisuju samo zaglavlje, kao u koraku 2b. Treća ispisuje
\code{working tree clean}: vježba nije mijenjala kod.

\textbf{Provjera ekstenzijom Database.} Pokušaj spajanja na vezu sada javlja
\emph{Can't connect to server on 127.0.0.1:5432} jer baza više ne radi.

\slika[0.5]{baza-nema-baze.png}{Database: veza bez baze}

Okruženje je spremno za sljedeće predavanje, koje počinje korakom 1.

\needspace{16\baselineskip}
\section{Sažetak naredbi}

Prvi terminal, korijen repozitorija:

\begin{kod}
git status                                  # working tree clean
git checkout predavanje-NN
cp .env.example .env                        # PowerShell: Copy-Item
docker ps -a                                # samo zaglavlje; inače docker rm -f -v
docker volume ls                            # samo zaglavlje; inače docker volume rm
docker compose up -d db
docker compose exec db psql -U pl_user -d pl_reg -c "\dt"
\end{kod}

Drugi terminal, mapa \code{api}:

\begin{kod}
.venv\Scripts\Activate.ps1                  # macOS/Linux: source .venv/bin/activate
python -c "import sys; print(sys.prefix)"   # ...\api\.venv
uv pip install -r requirements.txt
alembic upgrade head
alembic current                             # NNN (head)
python -m app.seed
uvicorn app.main:app --reload
\end{kod}

Na kraju: Ctrl+C, \code{deactivate}, pa u prvom terminalu \code{docker compose down -v}.

\section{Česti problemi}\label{sec:problemi}

Problemi s instalacijom alata opisani su u uputama \emph{Programiranje za web - priprema
razvojnog okruženja}. Ovdje su oni koji nastaju pri prelasku između predavanja.

\begin{center}\small
\begin{tabularx}{\linewidth}{P{4.6cm}P{3.8cm}L}
\arrayrulecolor{linije}\toprule
\zaglavlje \textbf{Simptom} & \textbf{Uzrok} & \textbf{Rješenje}\\
\midrule
\code{ModuleNotFoundError: No module named '...'} pri pokretanju \code{uvicorn} ili \code{alembic} &
  Novo predavanje traži paket koji nije instaliran &
  Korak 2c: aktivirajte okruženje i pokrenite \code{uv pip install -r requirements.txt}\\
\code{A virtual environment already exists at: .venv} &
  \code{uv venv} pokrenut u mapi koja već ima okruženje &
  Preskočite \code{uv venv}; za potpuno novo okruženje \code{uv venv --clear}, pa ponovno
  instalirajte pakete\\
\code{Can't locate revision identified by '...'} pri \code{alembic upgrade head} &
  Baza je na migraciji iz drugog predavanja, koju ovaj kod ne poznaje &
  \code{docker compose down -v}, pa korak 2b i 2d ispočetka\\
\code{alembic upgrade head} ne ispisuje nijedan \code{Running upgrade}, a \code{\textbackslash dt} prikazuje tablice &
  Baza iz prošlog predavanja nije obrisana & Isto kao redak iznad\\
\code{no configuration file provided: not found} &
  \code{docker compose} pokrenut iz mape \code{api} &
  Naredbu pokrenite u prvom terminalu, u korijenu repozitorija\\
\code{Bind for 0.0.0.0:5432 failed: port is already allocated} &
  Port 5432 drži drugi kontejner, npr. iz uputa za pripremu okruženja, ili lokalno
  instaliran PostgreSQL &
  Korak 2b: \code{docker ps -a}, pa \code{docker rm -f -v} za svaki kontejner; ako kontejnera
  nema, ugasite lokalni PostgreSQL\\
\code{Field required} za \code{DATABASE\_URL} ili drugu varijablu pri pokretanju &
  \code{.env} ne postoji ili u njoj nedostaje varijabla & Korak 2a\\
\code{error: pathspec 'predavanje-NN' did not match any file(s) known to git} &
  Lokalni repozitorij ne zna za novi tag & \code{git fetch --tags}, pa ponovite \code{git checkout}\\
\code{Your local changes ... would be overwritten by checkout} &
  Radno stablo nije čisto & Korak 1a\\
\code{ConnectionRefusedError} ili \code{Connection refused} pri \code{alembic} ili \code{python -m app.seed} &
  Baza nije pokrenuta &
  \code{docker compose ps}; ako nema reda s \code{Up}, \code{docker compose up -d db}\\
\bottomrule
\end{tabularx}
\end{center}

\end{document}
